\chapter{Cholecystitis (gallbladder inflammation)}
\index{Cholecystitis (gallbladder inflammation)}

\medskip
\noindent\textit{\textbf{Under review.} This chapter is an AI-assisted educational draft; it carries no source citations and is not a substitute for professional medical advice.}

\section*{Definition and Overview}
Cholecystitis is the medical term for inflammation of the gallbladder, the small pouch beneath the liver that stores the bile needed to digest fats. The inflammation is almost always triggered by a gallstone blocking the gallbladder's drainage tube, causing bile to become trapped and the gallbladder wall to swell and become infected. It typically comes on suddenly and produces severe pain in the upper right part of the belly. It is a common emergency and is usually cured by removing the gallbladder. There is also a chronic form, in which repeated milder irritation leaves the gallbladder scarred and thickened.

\section*{Epidemiology}
Gallbladder inflammation is one of the most frequent surgical emergencies in developed countries. It most often strikes adults over 40, and it is roughly twice as common in women as in men. Pregnancy, hormonal factors, obesity, and rapid weight loss all raise the risk because they encourage gallstone formation, and gallstones are the underlying cause of most attacks. In older people and those with diabetes, attacks tend to be more severe and more likely to progress to complications. In hospital intensive-care settings, a rarer form without stones (acalculous cholecystitis) can develop in the very ill.

\section*{Aetiology and Pathophysiology}
The chain of events usually begins with a gallstone, but any condition that blocks the gallbladder's drainage can cause it.

\begin{itemize}
    \item \textbf{Obstruction:} a gallstone lodges in the cystic duct, blocking the exit of the gallbladder
    \item \textbf{Stagnation:} trapped bile sits in the gallbladder, chemically irritating its lining
    \item \textbf{Inflammation:} the gallbladder wall becomes swollen, thickened, and painful
    \item \textbf{Infection:} bacteria move into the stagnant bile and multiply, adding infection on top of inflammation
    \item \textbf{Tense, stretched wall:} pressure builds inside the organ, reducing its blood supply and increasing the risk of tissue death
    \item \textbf{Chronic cholecystitis:} long-standing gallstones cause repeated low-grade irritation and scarring over years
\end{itemize}

\section*{Clinical Features}
The classic picture is a severe, steady pain that does not come and go like ordinary gallstone pain.

\begin{itemize}
    \item \textbf{Severe pain} in the upper right abdomen that may spread to the right shoulder or back
    \item \textbf{Pain made worse by breathing in deeply} or by pressure on the abdomen
    \item \textbf{Tenderness} just beneath the right ribs
    \item \textbf{Fever,} sometimes with chills
    \item \textbf{Nausea, vomiting, and loss of appetite}
    \item \textbf{Feeling generally unwell and lethargic}
    \item \textbf{Mild jaundice} in some cases if the main bile duct is affected too
\end{itemize}

\section*{Differential Diagnosis}
Because many conditions cause upper abdominal pain, the diagnosis is not always obvious from symptoms alone.

\begin{itemize}
    \item Biliary colic -- gallstone pain without inflammation
    \item Cholangitis -- infection of the bile ducts, often with jaundice and fever
    \item Acute pancreatitis
    \item Perforated or bleeding peptic ulcer
    \item Appendicitis, if the appendix sits high in the abdomen
    \item Kidney stones or infection of the right kidney
    \item Viral hepatitis or a liver abscess
    \item Pneumonia or pleurisy of the right lower lung
    \item Angina or heart attack presenting as upper abdominal pain
\end{itemize}

\section*{When to Seek Medical Care}
Abdominal pain that is severe, constant, or getting worse should always be assessed by a doctor, because cholecystitis can progress quickly and other dangerous conditions cause identical pain.

\textbf{Contact your local emergency services or attend an emergency department for:}
\begin{itemize}
    \item Severe pain that does not settle within a few hours
    \item Fever, chills, or rigors
    \item Yellow skin or eyes, dark urine, or pale stools
    \item Repeated vomiting
    \item A hard, swollen, or extremely tender abdomen
    \item Confusion, rapid breathing or heartbeat, or reduced urine output (possible sepsis)
\end{itemize}

\section*{Diagnosis and Investigations}
Investigations confirm the inflammation and check for gallstones and complications.

\begin{itemize}
    \item \textbf{History and physical examination:} the location and pattern of pain, and tenderness when pressing over the gallbladder
    \item \textbf{Blood tests:} elevated white cells and inflammatory markers (such as CRP), plus liver function tests
    \item \textbf{Abdominal ultrasound:} the first-choice scan; it shows gallstones, thickening of the gallbladder wall, and fluid around the organ
    \item \textbf{CT scan:} used when the ultrasound is unclear or a complication such as perforation is suspected
    \item \textbf{MRCP:} detailed imaging of the bile ducts if a duct stone is suspected
    \item \textbf{HIDA scan:} a nuclear medicine study that checks whether the gallbladder is draining normally
\end{itemize}

\section*{Management}
Cholecystitis is generally managed in hospital, and most people ultimately have the gallbladder removed.

\begin{itemize}
    \item \textbf{Fasting with intravenous fluids:} the bowel and gallbladder are rested while the attack is treated
    \item \textbf{Pain relief and anti-nausea medicines}
    \item \textbf{Antibiotics:} used when infection is present or suspected
    \item \textbf{Laparoscopic cholecystectomy:} keyhole removal of the gallbladder, performed early in the same admission in most cases
    \item \textbf{Open cholecystectomy:} needed occasionally when keyhole surgery is unsafe
    \item \textbf{Percutaneous drainage:} a tube through the skin to drain an infected gallbladder in people too unwell for surgery
    \item \textbf{ERCP:} to clear a stone from the main bile duct if present
\end{itemize}

\section*{Complications}
If left untreated, or in people with severe or delayed disease, serious complications can arise.

\begin{itemize}
    \item \textbf{Gangrene} of the gallbladder wall when its blood supply is cut off
    \item \textbf{Perforation:} the gallbladder ruptures and bile leaks into the abdomen
    \item \textbf{Abscess} around the gallbladder or liver
    \item \textbf{Bile peritonitis:} widespread inflammation of the abdominal lining
    \item \textbf{Cholangitis:} ascending infection of the bile ducts and bloodstream
    \item \textbf{Sepsis and septic shock}
    \item \textbf{Chronic cholecystitis} and recurrent attacks in those not treated surgically
\end{itemize}

\section*{Prognosis and Outcomes}
Most people with acute cholecystitis improve quickly with treatment, and the long-term outlook after the gallbladder is removed is excellent. After keyhole surgery most people leave hospital within a day or two and resume normal activities over a few weeks. Complications are uncommon in otherwise healthy people who are treated early, but the risk rises in older adults, people with diabetes, and those whose treatment is delayed. In people who are not fit for surgery, the inflammation often settles with antibiotics alone, but further attacks are likely, and some develop chronic symptoms.

\section*{Prevention and Self-Care}
\begin{itemize}
    \item Get abdominal pain assessed early rather than waiting for it to become severe
    \item Avoid rapid weight-loss diets and "crash" diets, which promote gallstone formation
    \item Keep to a healthy weight with a balanced, high-fibre diet and regular exercise
    \item Limit alcohol and stop smoking
    \item If you have known gallstones, report fever, jaundice, or persistent pain straight away
    \item After surgery, follow the discharge advice on diet, lifting, and activity, and attend follow-up
\end{itemize}
